The five isolated cubic pole certificates in WP26 of \cite{WeilCertificates}
are strengthened here to one continuous real interval. The
proof uses the actual finite-height theorem of Platt and Trudgian
and the unconditional zero-count bound of Hasanalizade, Shen
and Wong. Every unverified zero is retained in the tail bound;
no location of such a zero is assumed. The original factor
$h=3$, the cubic coefficient coordinates and both original
conductor norms are preserved by explicit receiving maps.

\section{The complete form and the two human theorems used}
For every real $\tau$ set
\begin{equation}
\omega_\tau=2+i\tau,\qquad h=3,\qquad
U_\tau(s)=\frac{s+1-i\tau}{2+i\tau-s},\qquad
\varphi_j^\tau(s)=\frac{\sqrt3}{2+i\tau-s}U_\tau(s)^j
\quad(0\le j\le3).
\tag{CP1}
\end{equation}
For $p=(p_0,p_1,p_2,p_3)^{\mathsf T}\in\C^4$ write
$F_p^\tau=\sum_{j=0}^3p_j\varphi_j^\tau$.
Use the full reflected Weil form
\begin{equation}
\W(F,G)=\sum_\rho m_\rho
       \overline{F(1-\bar\rho)}G(\rho),\qquad
T_3(\omega_\tau)=\bigl(\W(\varphi_i^\tau,\varphi_j^\tau)\bigr)_{i,j=0}^3 .
\tag{CP2}
\end{equation}
Here $\rho$ ranges over all distinct nontrivial zeta zeros,
and $m_\rho$ is its actual multiplicity. Absolute convergence
is proved below. Thus $p^*T_3p=\W(F_p^\tau,F_p^\tau)$.

The first human input is \cite[Theorem 1]{PlattTrudgian}, using only
the strictly smaller height
\begin{equation}
T_0=3\cdot10^{12}:
\qquad \Re\rho=\tfrac12
\quad\hbox{for every nontrivial zero with }|\Im\rho|\le T_0.
\tag{CP3}
\end{equation}
The published theorem reaches $3\,000\,175\,332\,800$.
The negative-height assertion in CP3 follows by conjugation.
The second input is \cite[Corollary 1.2]{HSW}, with its exact
original constants and counting convention:
\begin{equation}
N(T)=\sum_{0<\Im\rho\le T}m_\rho,\quad
|N(T)-f(T)|\le E(T)\quad(T\ge e),\qquad
f(T)=\frac{T}{2\pi}\log\frac{T}{2\pi e},\quad
E(T)=.1038\log T+.2573\log\log T+9.3675.
\tag{CP4}
\end{equation}
These are previously proved unconditional theorems, not
additional assumptions about the zeros being tested.

\section{Four separated actual zero intervals at every pole height}
For every $A\ge100$, CP4 implies
\begin{equation}
N(A+100)-N(A)>13\log A-52>0.
\tag{CP5}
\end{equation}
Here are all constants in the deduction. Since
$\log\log T\le\log T$ for $T\ge e$,
$E(T)\le\log T+10$. Also
$f'(T)=(2\pi)^{-1}\log(T/(2\pi))$ increases. Hence
\[
N(A+100)-N(A)
\ge\frac{100}{2\pi}\log\frac{A}{2\pi}
                 -\log A-\log(A+100)-20 .
\]
The elementary bounds $\pi<22/7$, $\log(2\pi)<2$ and
$\log2<1$ show that the right side is strictly larger than
$15(\log A-2)-2\log A-22=13\log A-52$.
Here $\log A-2>0$, so the inequality in its multiplier is
in the correct direction. Finally $\log100>4$:
$e<3$ gives $e^4<81<100$. For the other elementary bound,
$e^2>7>2\pi$ follows from the first five positive exponential
terms, whose sum is exactly seven. This proves CP5 including
its domain and endpoints.

Set
\begin{equation}
T_*=\frac{99}{100}T_0=2\,970\,000\,000\,000 .
\tag{CP6}
\end{equation}
For $|\tau|\le T_*$, let $a=|\tau|$ and choose the least
actual zero ordinate in each interval
\begin{equation}
I_j(a)=(a+100+200j,\ a+200+200j]\quad(j=0,1,2,3).
\tag{CP7}
\end{equation}
CP5 proves that each interval contains a zero. Such a least
ordinate exists because there are finitely many zeros in
each bounded interval. Every interval is strictly below
$T_0$, since $T_*+800<T_0$, so CP3 proves that these selected
zeros lie on the critical line. No simplicity is needed.
Write their positive ordinates as $\gamma_j$. For $\tau\ge0$
set $t_j=\gamma_j$; for $\tau<0$ set $t_j=-\gamma_j$.
These are four actual critical zeros in either case.
They are distinct, and their relative heights satisfy
\begin{equation}
100+200j<|t_j-\tau|\le200+200j\le800,\qquad
|t_j-t_i|\ge100(2|j-i|-1)\quad(i\ne j).
\tag{CP8}
\end{equation}
The intervals supply these inequalities directly; no continuity
of the selected zero labels as a function of $\tau$ is needed.
The resulting matrix bound will hold for every real $\tau$
in CP6, with the same numerical constant.

\section{An exact rational Vandermonde inverse bound}
For a critical zero of height $t$, put $z(t)=U_\tau(1/2+it)$.
Direct substitution into CP1 gives
\begin{equation}
z(t)=\frac{3/2+i(t-\tau)}{3/2-i(t-\tau)},\quad |z(t)|=1,\qquad
|z(t)-z(u)|
=\frac{3|t-u|}
 {\sqrt{9/4+(t-\tau)^2}\sqrt{9/4+(u-\tau)^2}}.
\tag{CP9}
\end{equation}
Retain the exact constants
\begin{equation}
L=800^2+\frac94=\frac{2560009}{4},\qquad
w_*=\frac3L,\qquad d_*=\frac{300}{L}.
\tag{CP10}
\end{equation}
For $z_j=z(t_j)$, each of the four row weights
$w_j=3/[9/4+(t_j-\tau)^2]$ is at least $w_*$, and
CP8--9 give
$|z_j-z_i|\ge d_*(2|j-i|-1)$.

Let $V=(z_j^\ell)_{0\le j,\ell\le3}$.
Its inverse has column $j$ equal to the coefficient vector of
the exact Lagrange polynomial
\begin{equation}
\ell_j(z)=\frac{\prod_{i\ne j}(z-z_i)}
                     {\prod_{i\ne j}(z_j-z_i)}.
\tag{CP11}
\end{equation}
Evaluation at the four distinct $z_i$ proves both inverse
compositions. Since $|z_i|=1$, the four numerator coefficient
moduli are at most $1,3,3,1$, so their squared Euclidean norm
is at most twenty. The three denominator separations are
at least $(d_*,3d_*,5d_*)$ for $j=0,3$ and
$(d_*,d_*,3d_*)$ for $j=1,2$. Therefore
\begin{equation}
\|V^{-1}\|_2^2\le\|V^{-1}\|_{\mathrm F}^2
\le20d_*^{-6}\left(\frac2{225}+\frac2{9}\right)
=\frac{208}{45}d_*^{-6}.
\tag{CP12}
\end{equation}
This bound retains all four separated intervals; it does not
replace their product by a determinant divided by a trace.

Taking one copy of each selected zero gives a positive Gram
$H_{\mathrm{sel},\tau}=V^*\diag(w_0,w_1,w_2,w_3)V$.
Because $p=V^{-1}Vp$, CP12 proves for every coefficient vector
\begin{equation}
p^*H_{\mathrm{sel},\tau}p\ge b_*\|p\|_2^2,\qquad
b_*=\frac{45}{208}w_*d_*^6
=\frac{135\cdot1200^6}{52\cdot2560009^7}>10^{-26}.
\tag{CP13}
\end{equation}
The last inequality is exact integer arithmetic: the difference
$135\cdot1200^6\cdot10^{26}-52\cdot2560009^7$ is the positive
integer
$2839912978528311935972297436813173952969845612$.
Multiplicities can only increase the contribution of these
verified zeros.

\section{Every unverified zero and the complete reflected tail}
Write a point in the closed strip as $s=\beta+it$,
$0\le\beta\le1$, and put $x=t-\tau$. For $|x|\ge10$,
the exact factors in CP1 imply
\[
|U_\tau(s)|^2
=\frac{(\beta+1)^2+x^2}{(2-\beta)^2+x^2}
\le1+\frac3{x^2}<\left(\frac{51}{50}\right)^2,\qquad
|U_\tau'(s)|=\frac3{|\omega_\tau-s|^2}\le\frac3{x^2}.
\]
Differentiate each of the four original functions without
changing its degree or the factor $\sqrt3$. For $0\le j\le3$,
\begin{equation}
|(\varphi_j^\tau)'(s)|
\le\frac{\sqrt3}{x^2}
\left[\left(\frac{51}{50}\right)^3+
 \frac9{|x|}\left(\frac{51}{50}\right)^2\right]
<\frac4{x^2},\qquad
|(F_p^\tau)'(s)|\le\frac8{x^2}\|p\|_2 .
\tag{CP14}
\end{equation}
For the strict numerical inequality it is enough to use
$\sqrt3<7/4$ and $|x|\ge10$, giving the rational upper
coefficient $3.495744<4$.
The last inequality is Cauchy--Schwarz on the four derivatives.

The reflected zeros $\rho=\beta+it$ and
$\rho^\vee=1-\bar\rho=1-\beta+it$ have equal multiplicity.
Integrating CP14 on their horizontal segment, whose length
is at most one and whose imaginary part stays $t$, gives
\[
|F_p^\tau(\rho)-F_p^\tau(\rho^\vee)|
\le\frac8{|t-\tau|^2}\|p\|_2.
\]
Every noncritical reflected pair of multiplicity $m$ consequently
has its exact contribution bounded below by
\begin{align}
2m\Re(\overline{F_p^\tau(\rho^\vee)}F_p^\tau(\rho))
&=m\bigl(|F_p^\tau(\rho)|^2+|F_p^\tau(\rho^\vee)|^2
 -|F_p^\tau(\rho)-F_p^\tau(\rho^\vee)|^2\bigr)\nonumber\\
&\ge-\frac{64m}{|t-\tau|^4}\|p\|_2^2.
\tag{CP15}
\end{align}
Critical zeros, whether verified or not, contribute nonnegatively.
For each positive-height noncritical pair, its two members
contribute $2m$ to $N$. The conjugate negative-height pair
has the same multiplicity and relative height $-t-\tau$.
Therefore the total possible negative contribution from all
unverified zeros is bounded by
\[
32\|p\|_2^2
\sum_{\Im\rho>T_0}m_\rho
\left[\frac1{|\Im\rho-\tau|^4}
      +\frac1{|\Im\rho+\tau|^4}\right].
\]
For $|\tau|\le T_*$ and $t>T_0$,
$|t\pm\tau|\ge t/100>T_0/100>10$. Thus CP14 applies,
and this last expression is at most
$64\cdot100^4\|p\|_2^2\sum_{\Im\rho>T_0}m_\rho/(\Im\rho)^4$.

The zero count CP4 also gives $N(T)\le T\log T$ for $T\ge100$.
In fact $f(T)<T\log T/6$, and
$E(T)\le\log T+10<T$, while $\log T>4$.
The last strict bound follows at $100$ and persists because
$T-\log T-10$ increases. Stieltjes integration by parts now gives
\begin{align}
\sum_{\Im\rho>T_0}\frac{m_\rho}{(\Im\rho)^4}
&=-\frac{N(T_0)}{T_0^4}
     +4\int_{T_0}^\infty\frac{N(t)}{t^5}\,dt\nonumber\\
&\le\frac4{T_0^3}
              \left(\frac{\log T_0}{3}+\frac19\right).
\tag{CP16}
\end{align}
The boundary term at infinity is zero by the same count bound.
Every term above the verified height has therefore been retained
in a convergent upper estimate.
For completeness CP4 also proves absolute convergence of CP2:
for fixed $\tau$, each $\varphi_j^\tau(\rho)=O(|\Im\rho|^{-1})$,
uniformly across the strip at large height, and
$\sum_{\Im\rho>R}m_\rho/(\Im\rho)^2<\infty$ follows by the
same Stieltjes calculation. Conjugate negative heights satisfy
the same bound. Hence the reflection grouping used above is valid.

Combining CP15--16, the whole unknown negative tail is at most
$e_*\|p\|_2^2$, where the following constants are exact:
\begin{equation}
e_*=\frac{256\cdot100^4}{T_0^3}\frac{91}{9}
=\frac{23296}{243\cdot10^{28}}
<\frac{24}{25}\,10^{-26},\qquad
\mu_*=b_*-e_*>4\cdot10^{-28}.
\tag{CP17}
\end{equation}
Here $\log T_0<30$. One elementary exact check is that the
finite sum $\sum_{k=0}^{30}30^k/k!$ exceeds $3\cdot10^{12}$,
so $e^{30}>T_0$. The rational comparison for $e_*$ is
$23296/243<96$, equivalent to $23296<23328$.
The last lower bound follows from CP13:
$10^{-26}-(24/25)10^{-26}=4\cdot10^{-28}$.

Every remaining zero below $T_0$ is critical by CP3 and
contributes nonnegatively. Adding CP13 and subtracting this
complete tail proves the continuous, unconditional result
\begin{equation}
\boxed{\quad
T_3(2+i\tau)\succeq\mu_* I_4
\succeq4\cdot10^{-28}I_4
\quad\hbox{for every real }|\tau|\le2\,970\,000\,000\,000.
\quad}
\tag{CP18}
\end{equation}
The same coefficient bound holds for all four coordinates
and all relative phases simultaneously. The selection of four
verified zeros proves one positive contribution; the rest
of the complete form is controlled, not omitted.

\section{Return to the unchanged original conductor and its Gamma norm}
Retain any fixed admitted reference conductor from FC21--32
of \cite{Fable029}, its actual symbol order $v$, its original shifts
and coefficients, and its exact source quotient and target:
\[
T_A:\mathcal P_{v+3}/\mathcal P_{v-1}\longrightarrow\mathcal P_3,
\qquad T_Ap(S')=\sum_{u,w=0}^8 a_{uw}p(S'+\beta_{8,u,w}).
\]
The reference roots are denoted $r_{\ell,*}$ and remain fixed
as the pole height $\tau$ varies. With the original marked
collision matrix $K$, let
\[
\ell_{\ell,*}(S)=
\prod_{b\ne\ell}\frac{S-r_{b,*}}{r_{\ell,*}-r_{b,*}},\qquad
y_\ell(S')=\sum_{u,w=0}^8a_{uw}
 (S'+\beta_{8,u,w})^v\ell_{\ell,*}(S'+\beta_{8,u,w}).
\]
Writing $Y$ for their columns in the original target monomials,
the retained map is $\Psi=YK^{-1}$. Its source partner is
$\Phi_{\rm FC}$ with $T_A\Phi_{\rm FC}=\Psi$.
All these original maps are the proved isomorphisms on the
actual quotient; the reference periods, moments and norms do
not depend on $\tau$.

Let $c_*=a_{\rm amp}/2-4$ and $\kappa=c_*-1/2$.
The translation matrix and rational numerator matrix are
\begin{align}
(C_\kappa)_{r\ell}&=\binom{\ell}{r}\kappa^{\ell-r}
 \quad(0\le r\le\ell\le3),\qquad
C_\kappa^{-1}=C_{-\kappa},\nonumber\\
R_\tau&=\sqrt3
\begin{pmatrix}
\omega^3&b\omega^2&b^2\omega&b^3\\
-3\omega^2&\omega^2-2b\omega&2b\omega-b^2&3b^2\\
3\omega&b-2\omega&\omega-2b&3b\\
-1&1&-1&1
\end{pmatrix},
\quad \omega=2+i\tau,\quad b=1-i\tau,\quad \det R_\tau=3^8.
\tag{CP19}
\end{align}
The $j$th column is exactly the coefficient vector of
$\sqrt3(z+b)^j(\omega-z)^{3-j}$.
For clarity its full inverse, also fixing every sign, is
\[
(R_\tau^{-1})_{j\ell}
=3^{-7/2}\sum_{r=0}^{\ell}
\binom{\ell}{r}\omega^r(-b)^{\ell-r}\binom{3-\ell}{j-r}.
\]
Indeed $u=(z+b)/(\omega-z)$ gives
$z=(\omega u-b)/(1+u)$ and $\omega-z=3/(1+u)$;
expanding $3^{-7/2}(1+u)^3p((\omega u-b)/(1+u))$
proves the inverse formula by substitution.
The determinant follows by changing to $t=\omega-z$:
the columns $\sqrt3(3-t)^jt^{3-j}$ become triangular
after reversal, and the two reversal signs cancel.

The exact original receiver and its inverse are consequently
\begin{equation}
B_\tau=R_\tau^{-1}C_\kappa\Psi,\qquad
B_\tau^{-1}=\Psi^{-1}C_{-\kappa}R_\tau,\qquad
F_{B_\tau q}^\tau(z)=\frac{(\Psi q)(z+\kappa)}{(\omega-z)^4}.
\tag{CP20}
\end{equation}
Multiplying the numerator coefficients proves the last identity.
It also proves the original conductor return
$B_\tau=R_\tau^{-1}C_\kappa T_A\Phi_{\rm FC}$.
Thus CP18 gives the actual, continuously parameterized form
on the original Fable coordinates:
\begin{equation}
\W_{\tau,\rm original}
=B_\tau^*T_3(2+i\tau)B_\tau
\succeq\mu_* B_\tau^*B_\tau
\quad(|\tau|\le T_*).
\tag{CP21}
\end{equation}
This is a strict positive form because $B_\tau$ is invertible.
The coefficient norm on the right is not substituted for
the original physical norm; its exact comparison is next.

For the original positive parameter $s_0\in\{1,a_{\rm amp}\}$,
retain the full target measure
\[
dm_{s_0}(y)=
\frac{(2\pi)^{s_0/2}2^{s_0/2}}{4\pi\Gamma(s_0/2)}
 \left|\Gamma(s_0/4+iy/2)\right|^2dy,
\quad
\|p\|_{\Gamma,W}^2=\int_{\R}|p(c_*+iy)|^2\,dm_{s_0}(y).
\]
Let $\mathcal H$ be this Gram in $(1,S',S'^2,S'^3)$ and
$G=\Psi^*\mathcal H\Psi$. In rational coefficients the same
norm has the exact Gram
\begin{equation}
H_\tau=B_\tau^{-*}GB_\tau^{-1}
=R_\tau^*C_{-\kappa}^*\mathcal HC_{-\kappa}R_\tau,\qquad
(H_\tau)_{ij}
=3\int_{\R}
 (3/2+i(y-\tau))^{3+j-i}
 (3/2-i(y-\tau))^{3+i-j}\,dm_{s_0}(y).
\tag{CP22}
\end{equation}
Here $H_\tau$ denotes the full native Gram; the selected-zero
Gram $H_{\mathrm{sel},\tau}$ was used only in CP13.
To prove CP22, evaluate its numerator at
$S'=c_*+iy$: the translated argument is exactly $1/2+iy$.
The two factors are conjugate, with exponents displayed
after multiplying conjugate column $i$ and column $j$.
This retains the original center, total mass and every cross term.

Put $M=(2\pi)^{s_0/2}$ and $\beta_0=s_0/2$. The original
Gamma moments are
$m_0=M$, $m_2=M\beta_0$,
$m_4=M(3\beta_0^2+2\beta_0)$,
$m_6=M(15\beta_0^3+30\beta_0^2+16\beta_0)$, with odd
moments zero, as proved for this unchanged measure in
FC56a/RW9 \cite{Fable029}. Hence every diagonal entry of CP22 is
the following explicit positive polynomial:
\begin{align}
g_0(\tau)=3\bigg[&
\frac{729}{64}m_0+\frac{243}{16}(m_2+\tau^2m_0)
+\frac{27}{4}(m_4+6\tau^2m_2+\tau^4m_0)\nonumber\\
&+m_6+15\tau^2m_4+15\tau^4m_2+\tau^6m_0\bigg],
\qquad \operatorname{tr}H_\tau=4g_0(\tau).
\tag{CP23}
\end{align}
Expansion of $3(9/4+(y-\tau)^2)^3$ proves this formula.
Since a positive Hermitian matrix has largest eigenvalue
at most its trace, CP18 yields the exact native lower bound
\begin{equation}
q^*\W_{\tau,\rm original}q
\ge\frac{\mu_*}{4g_0(\tau)}\,q^*Gq
\quad(|\tau|\le T_*).
\tag{CP24}
\end{equation}
Every coefficient of $g_0$ as a polynomial in $\tau^2$ is
positive. Thus $\mu_*/[4g_0(T_*)]$ is also a single explicit
positive lower constant for the full continuous interval
in the unchanged original target norm.

\section{All eight signed states and the exact remaining kernel}
In the original eight-state construction, retain the actual
four columns of pair means and half differences
$E=(e_1,\ldots,e_4)$, $O=(o_1,\ldots,o_4)$, with
$e_j=(q_{j,+}+q_{j,-})/2$ and
$o_j=(q_{j,+}-q_{j,-})/2$. These are the complete original
source points over the literal quartet, including both signs,
as defined and proved in SE1--7 of \cite{Fable029}. Their determinant
is
$1792\chi\delta^2\gamma^2(2\delta^2\gamma^2-\delta^2+\gamma^2)$,
where $\chi=\pm1$. It is nonzero on the original admitted
domain $0<\delta<1/2$, $\gamma>2$. Thus $O^{-1}$ is the
exact original inverse on this domain, with no extra state
or sign omitted. Write a full signed vector as $(\alpha,\beta)$,
whose actual evaluation is $z=E\alpha+O\beta$.

Let $G_Y=Y^*\langle\, ,\,\rangle_{\Gamma,W}Y$ be the
original target label Gram. The complete target norm and
arithmetic form are
\begin{align}
\mathbb H_W&=
\begin{pmatrix}I&0\\E&O\end{pmatrix}^{*}
\diag(G_Y,G)
\begin{pmatrix}I&0\\E&O\end{pmatrix},\nonumber\\
\mathbb W_\tau&=[E\ O]^*B_\tau^*T_3(2+i\tau)B_\tau[E\ O],
\qquad
S_\tau=\begin{pmatrix}I&0\\-O^{-1}E&O^{-1}B_\tau^{-1}\end{pmatrix}.
\tag{CP25}
\end{align}
Direct multiplication proves
\[
\begin{pmatrix}I&0\\E&O\end{pmatrix}S_\tau
=\diag(I,B_\tau^{-1}),\qquad
[E\ O]S_\tau=[0\ B_\tau^{-1}].
\]
Consequently CP18 proves the full simultaneous comparison
\begin{equation}
S_\tau^*\mathbb W_\tau S_\tau
=\diag(0_4,T_3(2+i\tau))
\succeq\diag(0_4,\mu_*I_4),\qquad
S_\tau^*\mathbb H_W S_\tau=\diag(G_Y,H_\tau).
\tag{CP26}
\end{equation}
The inverse of $S_\tau$ is
$\begin{psmallmatrix}I&0\\B_\tau E&B_\tau O\end{psmallmatrix}$,
verified by multiplying either way. The complete kernel, rank,
and target quotient coercivity are therefore exactly
\begin{equation}
\ker\mathbb W_\tau
=\{(\alpha,-O^{-1}E\alpha):\alpha\in\C^4\},\qquad
\operatorname{rank}\mathbb W_\tau=4,\qquad
\mathbb W_\tau|_{\rm quotient}
\succeq\frac{\mu_*}{4g_0(\tau)}
                         \mathbb H_W|_{\rm quotient}.
\tag{CP27}
\end{equation}
To specify the quotient norm without an informal deletion,
the full coordinate pair is $(\alpha,z)$ and its squared
norm is $\alpha^*G_Y\alpha+z^*Gz$.
For fixed $z$ the unique minimum over the kernel coset is
at $\alpha=0$. Its section is
$(0,O^{-1}z)$. In rational coordinates $d=B_\tau z$ it is
$(0,O^{-1}B_\tau^{-1}d)$ and has squared norm $d^*H_\tau d$.
This proves exactly the quotient and the constant in CP27.
There are no additional arithmetic null directions in this
continuous pole range, because CP18 makes $T_3$ positive definite.

The source norm remains distinct. Let
$G_X=X^*\langle\, ,\,\rangle_{\Gamma,U}X$ and
$G_U=\Phi_{\rm FC}^*\langle\, ,\,\rangle_{\Gamma,U}\Phi_{\rm FC}$,
with the original source quotient norm, and put
$H_{U,\tau}=B_\tau^{-*}G_UB_\tau^{-1}$.
For the actual source pair
$(X\alpha,\Phi_{\rm FC}z)$ the same multiplication gives
\begin{equation}
S_\tau^*\mathbb H_U S_\tau=\diag(G_X,H_{U,\tau}),\qquad
\mathbb W_\tau|_{\rm source\ quotient}
\succeq
\frac{\mu_*}{\operatorname{tr}H_{U,\tau}}\,
                 \mathbb H_U|_{\rm quotient}.
\tag{CP28}
\end{equation}
The source quotient again minimizes at $\alpha=0$;
$H_{U,\tau}$ is its exact positive Gram, with every original
inverse coefficient specified by FC24. The actual conductor
maps this pair to $(Y\alpha,\Psi z)$ because
$T_AX=Y$ and $T_A\Phi_{\rm FC}=\Psi$. Thus CP25--28 preserve
both original metrics and their precise morphism.

There is also a uniform source constant with every original
inverse weight explicitly present. Write the retained conductor
symbol as
$g_A(t)=t^{-v}e^{-4w(t)}E_A(w(t))=\sum_{n\ge0}c_nt^n$,
where $w(t)=-i\arctan t$ and
$c_0=(-i)^v\mu_v/v!\ne0$ is its actual leading moment.
The four inverse coefficients are exactly
\[
\eta_0=c_0^{-1},\quad \eta_1=-c_1/c_0^2,\quad
\eta_2=(c_1^2-c_0c_2)/c_0^3,\quad
\eta_3=(-c_1^3+2c_0c_1c_2-c_0^2c_3)/c_0^4 .
\]
Their convolution with $c_0,c_1,c_2,c_3$ gives $(1,0,0,0)$.
With the original weights $\rho_n=\sqrt{n!/(s_0/2)_n}$,
the actual inverse matrix in the original orthonormal source
and target frames has entry
$\eta_{j-i}\rho_j/\rho_{v+i}$ for $0\le i\le j\le3$,
and zero below the diagonal, by FC24/WCF6.
Thus its complete squared Frobenius norm and uniform
source comparison are
\begin{equation}
J_A=\sum_{0\le i\le j\le3}
 |\eta_{j-i}|^2\frac{\rho_j^2}{\rho_{v+i}^2}>0,\qquad
H_{U,\tau}\preceq J_AH_\tau,\qquad
\mathbb W_\tau|_{\rm source\ quotient}
\succeq\frac{\mu_*}{4J_Ag_0(T_*)}
                  \mathbb H_U|_{\rm quotient}
\quad(|\tau|\le T_*).
\tag{CP29}
\end{equation}
Indeed a matrix has operator norm at most its Frobenius norm,
so the literal source inverse gives
$\|T_A^{-1}p\|_{\Gamma,U}^2\le J_A\|p\|_{\Gamma,W}^2$.
Substitute $p=C_{-\kappa}R_\tau d$ to obtain the matrix
inequality; CP23--24 prove the last bound. Orthogonal
coordinates are used only to calculate the norm of this
unchanged physical inverse: all masses and quotient weights
are those of the original source and target.

These results establish a continuous positive cubic test
range and its exact full eight-state kernel. The proof uses
the original full-Weil sum and the original cubic receiver
throughout. It makes no assertion about degrees above three
or about poles outside the explicit interval CP6.

The certificate
\texttt{verify\_continuous\_pole\_cubic\_positivity.py}
checks all rational constant comparisons by integer arithmetic,
the coefficient-norm and reflected-pair multiplicity factors,
the exact chord identity, and the native diagonal-moment formula.
Its output is retained separately. The human theorems CP3--4
are cited mathematical inputs, not assertions recertified by
that short algebra script.
